\section{Introduction}

The rapid progress of large language models (LLMs), and in particular of
LLM-based agents, has attracted intense interest as a potential driver of
scientific discovery across a wide range of
domains~\cite{Lu2026AIScientist,Gottweis2026CoScientist,Mitchener2025Kosmos,Zhang2026AgenticMaterials}.
Modern LLMs combine broad scientific knowledge with flexible reasoning
capabilities, and when equipped with
tools~\cite{Yao2023ReAct,Anthropic2024MCP,Hou2025MCP},
skills~\cite{Wang2023Voyager,Jiang2026AgenticSkills},
memory~\cite{Chhikara2025Mem0,Gao2026SelfEvolvingAgents}, and reasoning strategies such
as planning and
reflection~\cite{Yao2023ReAct,Sun2023AdaPlanner,Shinn2023Reflexion},
the resulting
agents have the potential to drive scientific research flexibly and
autonomously.

To date, LLM agents have been developed toward the autonomous execution of
research in several fields.
Prominent examples include the end-to-end automation of artificial
intelligence research and engineering~\cite{Lu2026AIScientist,Jin2026Arbor}
and the automation of biomedical
discovery~\cite{Gottweis2026CoScientist,Ghareeb2025Robin,Mitchener2025Kosmos,Ke2025BioDisco},
as well as chemistry agents that couple LLMs with domain tools and robotic
laboratories~\cite{Bran2024ChemCrow,Boiko2023Coscientist,Darvish2025ORGANA}.
In materials science, LLM agents have been actively applied both to
computational materials
design~\cite{Jia2024LLMatDesign,Takahara2025MatAgent,Nduma2025Crystalyse,Kim2026Materealize,Zou2025ElAgente,Pham2026ChemGraph,Deng2026AtomisticSkills}
and to the automation of materials synthesis and
characterization~\cite{Fei2026AgenticLLM,YanguasGil2026ALDOptimization,shi2026knowledge,mandal2025evaluating,vriza2026operating},
and they are emerging as a practical route to accelerating the progress
of materials research~\cite{Zhang2026AgenticMaterials}.

The term ``autonomous research'' by LLM agents, however, encompasses
tasks with substantially different levels of
agency~\cite{RiosGarcia2026AIScientists}.
Much of the work above automates predefined workflows, in which the agent
reliably executes a prescribed sequence of computations, syntheses, or
measurements~\cite{Bran2024ChemCrow,Boiko2023Coscientist,Darvish2025ORGANA,Deng2026AtomisticSkills},
or performs optimization toward an externally specified
goal~\cite{YanguasGil2026ALDOptimization,Fei2026AgenticLLM}.
A task demanding substantially higher agency is to answer the question of
\emph{why}: to arrive at explanatory understanding that is earned through
repeated cycles of generating mechanistic hypotheses, devising the
empirical tests that can discriminate among them, and revising beliefs
against the resulting evidence.
Steps in this direction include agents that generate and rank scientific
hypotheses~\cite{Ghafarollahi2025SciAgents,Yang2024MOOSE,Yang2025MOOSEChem,Yang2025MOOSEChem2,Ke2025BioDisco,Gottweis2026CoScientist},
frameworks that place falsification at the center of the research
process~\cite{Liu2024AIGS,Huang2025POPPER}, belief-guided open-ended
exploration~\cite{Agarwal2025AutoDiscovery}, and closed-loop systems that
iteratively refine hypotheses against data, experiments, or
simulations~\cite{Mitchener2025Kosmos,Ghareeb2025Robin,Wang2026ScienceClaw,Jin2026Arbor,Wiemann2026DiscoverPhysics}.

Nevertheless, in existing agents, neither the agent itself nor the
human researcher can readily trace which hypotheses the agent
entertained, which tests it performed, and how the resulting evidence
moved its beliefs.
Hypotheses and beliefs typically remain implicit in free-text
reasoning or in undisclosed internal states.
Where the research record is externalized, what is tracked is
argument-based rankings~\cite{Gottweis2026CoScientist}, the provenance
of artifacts~\cite{Wang2026ScienceClaw}, or task-performance
scores~\cite{Jin2026Arbor}, rather than an explicit degree of belief
that empirical evidence can raise or lower.
Indeed, a recent large-scale analysis of agent execution traces found that
evidence was ignored in about two-thirds of traces and that beliefs were
revised after refutation in only about a quarter of
cases~\cite{RiosGarcia2026AIScientists}.
Without an explicit, auditable record of the
hypothesis--test--evidence--belief cycle, and without a structure that
disciplines how the cycle is operated, the scientific soundness of an
agent's conclusions cannot be
verified~\cite{Wang2026WorkflowClosure,Bisht2026AgenticNotBuilt}, and
the agent itself cannot systematically build on, revise, or retire its
own hypotheses over a long research horizon.

In this work, we propose the \emph{Hypothesis Evolution Protocol} (HEP),
an agent harness that turns the scientific procedure of proposing,
testing, and revising hypotheses in the light of evidence into an
explicit protocol, every step of which is traceable by both the LLM
agent and the human researcher.
HEP externalizes each hypothesis as a persistent object in an auditable
registry, where it carries a belief probability that can be moved
only by attached evidence, evolves through refinement and merging of parent
hypotheses, and progresses through a lifecycle from proposal to
resolution as supported, refuted, or dormant.
The protocol thus serves as a footing on which the agent operates the
hypothesis--test--evidence--belief cycle in a disciplined way.
Using materials-science research tasks as a testbed, we show that an agent
equipped with HEP actually operates the hypothesis--test--evidence--belief
cycle that planning-style agents lack, that the same protocol generalizes
across research questions, and that the depth
to which the protocol is exploited scales with base-LLM capability.
